% diagrams/firstwords/tesoro.tex -- el viernes del cumpleaños de Dani: al
% final de la búsqueda, ¡el tesoro! Un cofre lleno de monedas de
% chocolate, Dani con los brazos arriba y Toby, que quiere probarlas.
\begin{tikzpicture}[dibujofw]
  \begin{scope}[shift={(-4.6,0)}, scale=0.95]
    \ninoFW{Dani}{Arriba}{Arriba}
    \begin{scope}[shift={(0,6.3)}] \bocaOFW \end{scope}
  \end{scope}
  % la tapa, abierta
  \filldraw[fill=white] (-2.4,2.4) -- (-2.1,4.6) .. controls (-1.0,5.3) and (1.0,5.3) .. (2.1,4.6) -- (2.4,2.4) -- cycle;
  \draw (-2.25,3.5) .. controls (-1.0,4.1) and (1.0,4.1) .. (2.25,3.5);
  % las monedas
  \foreach \x/\y in {-1.6/2.55, -0.55/2.6, 0.55/2.6, 1.6/2.55, -1.1/3.2, 0/3.3, 1.1/3.2}
    {\filldraw[fill=white] (\x,\y) circle (0.52); \draw (\x,\y) circle (0.3);}
  % el cofre, con sus bandas y la cerradura
  \filldraw[fill=white] (-2.6,0) rectangle (2.6,2.5);
  \filldraw[fill=white] (-2.0,0) rectangle (-1.5,2.5); \filldraw[fill=white] (1.5,0) rectangle (2.0,2.5);
  \filldraw[fill=white] (-0.4,1.1) rectangle (0.4,1.9);
  \fill (0,1.6) circle (0.1); \draw (0,1.6) -- (0,1.3);
  % monedas por el suelo
  \foreach \x in {-3.1,3.0} {\filldraw[fill=white] (\x,0.18) ellipse (0.42 and 0.18);}
  \begin{scope}[shift={(4.6,0)}, scale=0.75] \tobySentadoFW \end{scope}
  \draw (-7.0,0) -- (7.0,0);
\end{tikzpicture}
